\section{Síntesis y ampliación}

La semana 8 mostró cómo la IA transforma la construcción de aplicaciones: de la pila LAMP tradicional al paradigma impulsado por IA de lenguaje natural $\rightarrow$ aplicación. El análisis detallado de la arquitectura de v0 reveló que detrás de la «magia» del App Builder hay un workflow y un composite model cuidadosamente diseñados.

\subsection{Puntos clave}

\begin{itemize}
    \item El desarrollo web evolucionó de LAMP $\rightarrow$ MERN $\rightarrow$ JAM $\rightarrow$ Serverless $\rightarrow$ AI App Builder
    \item La IA hace que los ingenieros sean más multifuncionales, y las personas sin conocimientos técnicos también pueden construir aplicaciones
    \item Arquitectura de v0: Intent $\rightarrow$ Context $\rightarrow$ Code Gen $\rightarrow$ Validation $\rightarrow$ Human Review
    \item «Agent es Workflow» — lo esencial no es una única llamada al modelo, sino una cadena de procesamiento de varios pasos
    \item Stream manipulation y composite model family resuelven las limitaciones de los modelos de propósito general
    \item La mayor limitación: cuando ocurre un error, todavía se necesita que una persona entienda el código subyacente para depurarlo
\end{itemize}

\subsection{Hoja de ruta de implementación}

\begin{table}[H]
\centering
\begin{tabular}{p{0.18\textwidth} p{0.26\textwidth} p{0.3\textwidth}}
\toprule
Etapa & Objetivo principal & Práctica típica \\
\midrule
Validación de prototipo & Validar rápidamente si la idea es viable & Generar con lenguaje natural páginas, formularios y el esqueleto de las interacciones \\
Integración con el sistema & Conectar datos reales y el sistema de diseño & Conectar APIs, completar tipos y alinear las especificaciones de los componentes \\
Puesta en producción & Garantizar un mantenimiento estable y la colaboración de varias personas & Agregar pruebas, hacer review, y consolidar el workflow y los mecanismos de reversión \\
\bottomrule
\end{tabular}
\caption{Ruta típica de adopción de AI App Builder}
\end{table}

\subsection{Lista de acciones para el equipo}

\begin{enumerate}
    \item Elegir primero una herramienta interna de bajo riesgo o un proyecto de prototipo para hacer un piloto, en lugar de aplicarlo directamente a las páginas centrales del negocio
    \item Definir de antemano el sistema de diseño y los límites de los componentes, para que el resultado generado se ajuste de manera estable a las normas de ingeniería existentes
    \item Incorporar el resultado generado por la IA al proceso normal de code review, pruebas y reversión, en lugar de tratarlo como un proceso excepcional
\end{enumerate}

\subsection{Lecturas adicionales}

\begin{itemize}
    \item Vercel v0: \url{https://v0.dev} (gratis para estudiantes: \url{https://v0.app/students})
    \item Vercel AI SDK: \url{https://ai-sdk.dev/docs/agents/overview}
    \item Publicación de blog sobre Composite Model Family: \url{https://vercel.com/blog/v0-composite-model-family}
    \item Lovable: \url{https://lovable.dev}
    \item Bolt.new (StackBlitz): \url{https://bolt.new}
    \item Replit: \url{https://replit.com}
\end{itemize}

%% --- Fin del contenido --- %%

\end{document}
